\section{Trained Ctl-1: learned residual geometry}

The analytic task fixes the mechanism by construction. We next ask whether the
same geometry appears in a learned residual space. We train a small Transformer
on sequences with two latent predicates and an interactional target. We extract
the final residual representation, fit probes for $x_1$, $x_2$, and
$x_1x_2$, and build first-order and lifted estimator--direction pairs.
Directions are normalized to the same initial target-head gain.

This experiment is a sanity check. The auxiliary feature heads are supervised,
so the experiment shows that the geometry survives in a learned residual
representation; it does not show that a Transformer discovers the interaction
without supervision.

The null again behaves correctly. At $\gamma=0$, first-order and lifted pairs
match on target control and collateral. As $\gamma$ grows, first order remains
useful for the target but tends to pay more collateral at the default nuisance
placement. The numerical separation is smaller and more variable than in the
analytic task. Some collateral ratios are denominator-sensitive because lifted
collateral is close to zero, so we use the sweeps as geometry diagnostics rather
than as headline effect sizes.

\begin{figure}[t]
  \centering
  \begin{subfigure}[t]{0.49\linewidth}
    \centering
    \includegraphics[width=\linewidth]{figure_03_trained_ctl1_interaction_sweep.png}
    \caption{Interaction strength.}
  \end{subfigure}
  \hfill
  \begin{subfigure}[t]{0.49\linewidth}
    \centering
    \includegraphics[width=\linewidth]{figure_04_trained_ctl1_nuisance_sweep.png}
    \caption{Nuisance placement.}
  \end{subfigure}
  \caption{Trained Ctl-1. The learned task preserves the qualitative geometry,
  while the intervals expose denominator sensitivity and training variation.}
  \label{fig:trained-ctl1}
\end{figure}
